\documentclass[11pt,a4paper]{article}
\usepackage[margin=1in]{geometry}
\usepackage{hyperref}
\usepackage{enumitem}
\usepackage{xcolor}
\usepackage{listings}

\lstset{
  basicstyle=\ttfamily\small,
  breaklines=true,
  frame=single,
  columns=fullflexible,
}

\title{\bfseries Week 2: Profile Page, From Mock Data to the
  Real Database}
\author{Saksham Boran \\ \small Roll No: 1024030856}
\date{8 September 2026}

\begin{document}
\maketitle

\section{Context}
Google Sign-In (college-domain-restricted, via Auth.js v5 with
database sessions) and the onboarding flow landed early this week,
which meant a signed-in, onboarded user with real \texttt{Profile}
and \texttt{EvidenceRecord} rows in the database finally existed.
That unblocked the task queued up from Week 1: a profile page for
that user to view their own data.

\section{Frontend: Profile Page, First Pass}
Built \texttt{app/profile/page.tsx} against mock data first,
deliberately shaped to mirror the two onboarding-path variants
already defined in the schema (Path A: resume/GitHub evidence;
Path B: self-reported comfort level, interests, availability), so
the layout would be validated against the real data shapes before
wiring up any queries. Also added a ``View profile'' link next to
``Sign out'' on the signed-in home card, as the only entry point to
the new page.

\section{Backend: From Mock Data to a Real Query}
Once the shape held up, relocated the page into the
\texttt{(protected)} route group -- which gates on completed
onboarding at the layout level, using the session from the Google
sign-in flow built earlier in the week -- and replaced the mock
fixtures with a real query:
\begin{itemize}[leftmargin=0.7cm]
  \item Profile fields read from \texttt{prisma.profile}.
  \item Identity (name, email, avatar) read from the session
    object rather than duplicated into \texttt{Profile}.
  \item Resume/GitHub connection status derived from the user's
    \texttt{EvidenceRecord} rows (presence/absence and
    \texttt{source} field), rather than a separate boolean flag
    that could drift out of sync with the actual evidence.
\end{itemize}
The route group's layout already redirects a not-yet-onboarded
user away, so the page itself only needed a direct-navigation guard
for someone hitting the URL manually mid-onboarding -- the common
path was already covered upstream.

\section{Debugging: Broken Avatars With No Fallback}
\textbf{Error encountered:} the profile header rendered a plain
\texttt{<img>} whenever \texttt{Profile.image} was non-null. This
showed a broken-image icon instead of a sensible fallback, because
Google's avatar URLs (\texttt{lh3.googleusercontent.com}) both
rate-limit by IP and expire.

\textbf{Key observation:} the existing initials-fallback logic
only covered the \emph{null} case (no image URL at all) --  a URL
that resolves at render time but fails to \emph{load} in the
browser is a different failure mode that needs handling on the
client, not the server.

\textbf{Solution:} extracted a client component
(\texttt{Avatar.tsx}) that swaps to initials on the image's
\texttt{onError} event, plus a mount-time check for failures that
happen before hydration attaches that handler (the image is
server-rendered, so the browser can fail to load it before any
client JS has run):
\begin{lstlisting}[language=JavaScript]
function Avatar({ src, initials }) {
  const [failed, setFailed] = useState(false);
  useEffect(() => {
    // catches a failure that happened before hydration
    if (imgRef.current && !imgRef.current.complete) setFailed(true);
  }, []);
  if (failed || !src) return <InitialsFallback text={initials} />;
  return <img ref={imgRef} src={src} onError={() => setFailed(true)} />;
}
\end{lstlisting}

\section{Verification}
No automated test framework exists yet, so verification was
manual: a live walkthrough against the real Neon database and a
real Google account confirmed the profile page correctly renders
both onboarding-path variants, resume/GitHub status reflects the
actual \texttt{EvidenceRecord} rows (not a stale flag), and a
simulated broken avatar URL falls back to initials instead of a
broken-image icon.

\section{Outcome}
Shipped the profile page as three small, separately reviewed
commits: mock-data version, database-backed version, then the
avatar fix -- each landing cleanly rather than as one large
first-pass diff.

\section{Follow-ups}
\begin{itemize}[leftmargin=0.7cm]
  \item At this point the only evidence source was an unverified
    GitHub URL a user typed in during onboarding -- the profile
    page displayed it as-is. Turning that into a real, verified
    GitHub connection was the natural next task, picked up in
    Week 3.
  \item No automated tests cover the profile page yet; correctness
    still relies on manual walkthroughs.
\end{itemize}

\end{document}
